\section*{Bab \ref{ch:b3:total-synthesis} --- Sintesis Total: Strategi dan Rute Klasik}

\begin{solution}{exo:b3:total-synthesis:1}
Pada 90\,\%: $0.9^{10} = 35$\,\% dan $0.9^{20} = 12$\,\%. Pada 80\,\%: $0.8^{10} = 11$\,\% dan $0.8^{20} = 1.2$\,\%.
\end{solution}

\begin{solution}{exo:b3:total-synthesis:2}
Seluruhnya $5 + 3 + 1 + 4 = 13$ tahap; LLS $5 + 1 + 4 = 10$.
\end{solution}

\begin{solution}{exo:b3:total-synthesis:3}
4-(Dimetilamino)butan-2-on, $\ce{CH3COCH2CH2N(CH3)2}$: enol aseton teradisi pada ion
iminium $\ce{CH2=N+(CH3)2}$.
\end{solution}

\begin{solution}{exo:b3:total-synthesis:4}
Diskoneksikan kedua ikatan cincin yang alilik terhadap C=C: buta-1,3-diena dan but-3-en-2-on
(metil vinil keton), yang bereaksi jika dipanaskan.
\end{solution}

\begin{solution}{exo:b3:total-synthesis:5}
Linear: $0.9^{21} = 11$\,\%. Konvergen: $0.9^{11} = 31$\,\%, hampir tiga kali lebih banyak.
\end{solution}

\begin{solution}{exo:b3:total-synthesis:6}
A: $(7 + 1)/12 = 67$\,\%. B: $7/9 = 78$\,\%.
\end{solution}

\begin{solution}{exo:b3:total-synthesis:7}
Tanpa gugus pelindung: $0.9^8 = 43$\,\%. Dengan empat tahap tambahan pada 95\,\%: $43\,\% \times 0.95^4 = 35$\,\%.
\end{solution}

\begin{solution}{exo:b3:total-synthesis:8}
Primer: eter silil yang meruah (\textit{tert}-butildifenilsilil), dilepaskan dengan fluorida.
Sekunder: eter 4-metoksibenzil, dilepaskan melalui oksidasi dengan DDQ. Alkohol
tersier, yang terhalang, sering dapat dibiarkan bebas, atau dilindungi sebagai eter silil kecil
yang diputus oleh asam lemah. Setiap kondisi tidak menyentuh gugus-gugus lainnya.
\end{solution}

\begin{solution}{exo:b3:total-synthesis:9}
Sintesis-sintesis total itu memerlukan puluhan tahap dengan \omterm{def:b3:total-synthesis:architecture}{rendemen keseluruhan} jauh di bawah 1\,\%;
senyawa yang memiliki sistem cincin lengkap dan sebagian besar pusat stereo, yang diekstraksi dari
daun jarum yew Eropa, sebuah sumber yang dapat diperbarui, diubah dalam beberapa tahap saja.
\end{solution}

\begin{solution}{exo:b3:total-synthesis:10}
Adisi Michael: enolat dion itu (C2, di antara kedua karbonil) teradisi pada
karbon-$\beta$ metil vinil keton, sehingga terbentuk triketon. Aldol: enolat metil keton
menyerang sebuah karbonil cincin secara intramolekular, dan dehidrasi memberikan
sikloheksenon: keton Wieland--Miescher, sebuah enadion bisiklik dengan gugus metil
angular.
\end{solution}

\begin{solution}{exo:b3:total-synthesis:11}
$y^{m+1}/y^{2m+1} = y^{-m}$: pada 90\,\%, 1.7 untuk $m = 5$ dan 4.9 untuk $m = 15$. Makin panjang rutenya, makin
menguntungkan konvergensi.
\end{solution}

\begin{solution}{exo:b3:total-synthesis:12}
Setiap cincin menutup dengan kation dan ikatan rangkap penyerang dalam susunan
mirip kursi; adisi pada setiap ikatan rangkap berlangsung anti, dan substituen-substituennya menempati
posisi pseudo-ekuatorial, yang membuat setiap fusi cincin \textit{trans}. Ikatan rangkap \textit{E}
menempatkan kedua lanjutan rantainya sedemikian rupa sehingga fusi \textit{trans} terbentuk; ikatan rangkap \textit{Z} akan
memberikan fusi \textit{cis}: geometri poliena menyandikan stereokimia
produknya.
\end{solution}

\begin{solution}{pb:b3:total-synthesis:1}
\textbf{1.} $\ce{C4H6O2}$, $\ce{CH5N}$, $\ce{C5H6O5}$, $\ce{C8H13NO}$.

\textbf{2.} $\ce{C4H6O2 + CH5N + C5H6O5 -> C8H13NO + 2CO2 + 2H2O}$.

\textbf{3.} 86.0, 31.0, dan \qty{146.0}{g/mol}; tropinon \qty{139.0}{g/mol}.

\textbf{4.} $139.0/(86.0 + 31.0 + 146.0) = 0.53$: 53\,\%.

\textbf{5.} Gugus-gugus $\ce{CH2}$ tengahnya, yang masing-masing berada di antara dua gugus karbonil, mudah mengenol di dalam
air; aseton sendiri terlalu lemah nukleofilitasnya, dan gugus-gugus karboksil pengaktif
lepas sesudahnya.

\textbf{6.} Ion-ion iminium memerlukan amina bebas (asam yang terlalu banyak memprotonasinya), sedangkan
pembentukan enol dan iminium memerlukan sedikit asam: pH di dekat netral menyeimbangkan keduanya.

\textbf{7.} $\ce{CH3NH2}$ teradisi pada satu gugus aldehida; lepasnya air memberikan
$\ce{R-CH=N+H-CH3}$, sebuah ion iminium.

\textbf{8.} Sebuah karbon enol asam asetondikarboksilat menyerang karbon iminium, sehingga membuat
ikatan C--C dan sebuah amina sekunder.

\textbf{9.} Amina sekunder itu berkondensasi dengan gugus aldehida kedua pada rantai
yang sama, sehingga terbentuk ion iminium siklik (beranggota lima).

\textbf{10.} Enol di sisi lain keton menyerangnya, sehingga menutup
cincin beranggota enam dan kerangka bisiklik.

\textbf{11.} Masing-masing adalah asam $\beta$-keto, yang melepaskan $\ce{CO2}$ melalui keadaan transisi siklik beranggota enam,
sehingga terbentuk enol keton itu.

\textbf{12.} Dua ikatan C--C dan dua ikatan C--N.

\textbf{13.} $0.75^{14} = 0.018$: 1.8\,\%.

\textbf{14.} $42/1.78 = 24$.

\textbf{15.} 100\,\%: satu tahap, seluruhnya konstruksi ikatan.

\textbf{16.} Satu operasi, air sebagai pelarut, komponen-komponen yang murah, tanpa gugus pelindung, dan
hampir tanpa limbah selain karbon dioksida dan air.

\textbf{17.} Tumbuhan membangun kerangka tropana dari ion iminium siklik yang berasal dari asam amino
dan sebuah satuan yang berasal dari asetat, melalui kimia Mannich
yang sejenis.

\textbf{18.} Kaskade (dan \omterm{def:b3:pericyclic:families}{sikloadisi} pembentuk cincin seperti reaksi
Diels--Alder).

\textbf{19.} Tidak: sebuah \omterm{def:b3:point-groups:mirror}{bidang cermin} melalui N, gugus metilnya, C3, dan oksigen
karbonil.

\textbf{20.} Masing-masing memiliki empat gugus yang berbeda, tetapi keduanya saling bayangan cermin terhadap
bidang simetri molekul itu sendiri: senyawa bertipe meso.

\textbf{21.} Diastereomer, yang berbeda dalam orientasi OH pada C3 relatif terhadap
jembatan nitrogen.

\textbf{22.} Tidak: \omterm{def:b3:point-groups:mirror}{bidang cermin} dipertahankan pada keduanya.

\textbf{23.} Kedua muka gugus karbonil tidak setara: yang satu menghadap jembatan
nitrogen, yang lain menghadap jembatan dua karbon, dan keduanya menghalangi secara berbeda.

\textbf{24.} Rute satu wadah memberikan sekitar 24 kali \omterm{def:b3:total-synthesis:architecture}{rendemen keseluruhan} rute linear
14 tahap.
\end{solution}
